\DocumentMetadata{
  pdfversion=2.0,
  pdfstandard=ua-2,
  lang=en-US,
  tagging=on,
  tagging-setup={math/setup=mathml-SE,tabsorder=structure}
}
\documentclass[11pt,letterpaper]{article}
\usepackage[margin=0.68in,includefoot]{geometry}
\usepackage{newcomputermodern}
\usepackage{amsmath,amscd,mathtools}
\usepackage{tikz,adjustbox}
\providecommand{\square}{\mdlgwhtsquare}
\usepackage[hidelinks]{hyperref}
\hypersetup{
  pdftitle={Algebra First-Year Exam, May 2025, Part 1},
  pdfauthor={University of Florida Department of Mathematics},
  pdfsubject={Graduate mathematics examination},
  pdfdisplaydoctitle=true,
  bookmarksopen=true
}
\setcounter{secnumdepth}{0}
\setlength{\parindent}{0pt}
\setlength{\parskip}{0.28em}
\setlength{\emergencystretch}{3em}
\setlength{\leftmargini}{2.15em}
\setlength{\leftmarginii}{2.65em}
\setlength{\leftmarginiii}{3em}
\pagestyle{empty}
\raggedbottom
\allowdisplaybreaks
\NewTaggingSocketPlug{block/list/label}{examproblem}{%
  \tagstructbegin{tag=Lbl}\tagstructend%
  \tagstructbegin{tag=\LBody}%
  \tagstructbegin{tag=\ProblemHeadingTag,title-o={\ProblemHeadingText}}%
  \tagmcbegin{tag=\ProblemHeadingTag,actualtext={\ProblemHeadingText}}%
  #2%
  \tagmcend%
  \tagstructend%
}
\newenvironment{examproblems}[2]{%
  \def\ProblemHeadingTag{#2}%
  \AssignTaggingSocketPlug{block/list/label}{examproblem}%
  \begin{enumerate}%
  \setcounter{enumi}{#1}%
  \renewcommand{\labelenumi}{\textbf{\theenumi.}}%
  \setlength{\topsep}{0.35em}%
  \setlength{\partopsep}{0pt}%
  \setlength{\itemsep}{0.48em}%
  \setlength{\parsep}{0.18em}%
}{\end{enumerate}}
\newenvironment{examsubparts}{%
  \AssignTaggingSocketPlug{block/list/label}{default}%
  \begin{enumerate}%
  \setlength{\topsep}{0.18em}%
  \setlength{\partopsep}{0pt}%
  \setlength{\itemsep}{0.16em}%
  \setlength{\parsep}{0.1em}%
}{\end{enumerate}}
\newenvironment{examinstructions}{%
  \par\begingroup\itshape
}{%
  \par\endgroup\vspace{0.18em}
}
\makeatletter
\renewcommand\subsection{\@startsection{subsection}{2}{0pt}%
  {-1.25ex plus -.25ex minus -.1ex}%
  {0.55ex plus .1ex}%
  {\normalfont\large\bfseries}}
\renewcommand{\@maketitle}{%
  \newpage\null\vskip -1em%
  {\footnotesize\bfseries
    \href{https://www.ufl.edu/}{University of Florida}, \href{https://math.ufl.edu/}{Department of Mathematics}\par}%
  \vskip -0.5em%
  \tagstructbegin{tag=H1,title={Algebra First-Year Exam, May 2025, Part 1}}%
  \tagpdfparaOff%
  \tagmcbegin{tag=H1}%
    {\LARGE\bfseries\href{https://gma.math.ufl.edu/exams/algebra-first-year/}{Algebra First-Year Exam}, May 2025, Part 1\par}%
  \tagmcend%
  \tagpdfparaOn%
  \tagstructend%
  \vskip 0.25em%
  \tagmcbegin{artifact}%
    \hrule height 1pt%
  \tagmcend%
  \vskip 1em%
}
\makeatother
\title{Algebra First-Year Exam, May 2025, Part 1}
\author{}
\date{}
\begin{document}
\maketitle
\thispagestyle{empty}
\begin{examinstructions}
Answer 4 questions. Indicate which 4 problems you wish to be graded. Write your solutions clearly in complete English sentences. You may quote theorems (within reason) as long as you state them clearly.
\end{examinstructions}
\begin{examproblems}{0}{H2}
\def\ProblemHeadingText{Problem 1.}
\item
Let~\(p\) be a prime and~\(G\) be a nontrivial finite~\(p\)-group.
\begin{examsubparts}
\item[\textnormal{(a)}]
Prove that \(Z(G)\) is not trivial.
\item[\textnormal{(b)}]
Prove that every nontrivial finite~\(p\)-group has a normal subgroup of index~\(p\).
\end{examsubparts}
\def\ProblemHeadingText{Problem 2.}
\item
Let \(\sigma\in S_n\) be an~\(n\)-cycle.
\begin{examsubparts}
\item[\textnormal{(a)}]
Prove that \(\sigma\in A_n\) if and only if~\(n\) is odd.
\item[\textnormal{(b)}]
Prove that \(C_{S_n}(\sigma)=\langle\sigma\rangle\).
\item[\textnormal{(c)}]
Prove that when~\(n\) is odd, \(A_n\) has two conjugacy classes of~\(n\)-cycles.
\end{examsubparts}
\def\ProblemHeadingText{Problem 3.}
\item
Let~\(G\) be a group and let \(T(G)=\{x\in G:\exists n\geq 1, x^n=1\}\).
\begin{examsubparts}
\item[\textnormal{(a)}]
Prove that if~\(G\) is abelian then \(T(G)\) is a subgroup of~\(G\).
\item[\textnormal{(b)}]
Give an example of a nonabelian group~\(G\) such that \(T(G)\) is not a subgroup of~\(G\). Justify any claims that you make.
\end{examsubparts}
\def\ProblemHeadingText{Problem 4.}
\item
This problem has two parts.
\begin{examsubparts}
\item[\textnormal{(a)}]
Find the order of the group \(\operatorname{GL}(4,\mathbb{Z}/3\mathbb{Z})\).
\item[\textnormal{(b)}]
Show that there exists a nonabelian group of order \(2025=3^4\cdot 5^2\) in which the Sylow~\(3\)-subgroup is abelian.
\end{examsubparts}
\def\ProblemHeadingText{Problem 5.}
\item
Let~\(K\) and~\(H\) be groups, with~\(K\) abelian. Consider a homomorphism \(\phi:K\to\operatorname{Aut}(H)\) and the associated semidirect product \(G=H\rtimes_\phi K\).
\begin{examsubparts}
\item[\textnormal{(a)}]
Prove that when~\(K\) is identified with the subgroup \(\{(1,k)\mid k\in K\}\) of~\(G\), \(\operatorname{Ker}(\phi)\) is contained in the center of~\(G\).
\item[\textnormal{(b)}]
Prove that every group of order \(44\) has a cyclic subgroup of order \(22\).
\end{examsubparts}
\def\ProblemHeadingText{Problem 6.}
\item
Let~\(G\) be a finite group, and let~\(P\) be a Sylow~\(p\)-subgroup of~\(G\) for some prime~\(p\). Let \(N=N_G(P)\), the normalizer of~\(P\). Prove that \(N_G(N)=N\).
\end{examproblems}
\end{document}
